\documentclass[a4paper,11pt]{article}
\usepackage{exam-style-842-unit}
\examinehead{842 信息化概论 · 第 05 单元练习册}{练习二\quad 云服务与虚拟化}
\begin{document}

\unitcover{842 信息化概论 · 单元练习册}{第 05 单元\quad 云计算与大数据}{练习二\quad 云服务与虚拟化（共 100 分）}{约 80 分钟}{IaaS、PaaS、SaaS 的分层辨析与典型例子；虚拟化技术（服务器／存储／网络虚拟化、Hypervisor、容器与虚拟机）；主流云平台（AWS、阿里云等）；云安全与数据隐私。}

\begin{knowmap}
\textbf{三层模式：}按 NIST 划分 IaaS（第 1 层，最底层）$\rightarrow$ PaaS（第 2 层）$\rightarrow$ SaaS（第 3 层）；自下而上抽象程度提高，用户直接管理的范围递减\\[2pt]
\textbf{IaaS：}向用户提供服务器、存储、网络等基本计算资源；用户可自由安装和定制操作系统，只需专注自己的应用；教材典型服务为 EC2、EBS\\[2pt]
\textbf{PaaS：}向用户提供分布式操作系统、分布式数据库等基础软件，公司开发在这一层进行；亚马逊因创立 PaaS 商业模式被评为 2008 年创新企业\\[2pt]
\textbf{SaaS：}用户通过网页浏览器或智能终端使用远程服务器上的应用；微信、微博、云存储、在线视频、Google Apps、Dropbox、iCloud 均属此类\\[2pt]
\textbf{虚拟化：}在软件中仿真计算机硬件，以虚拟资源提供服务；教材列“一虚多”“多虚一”两种模式；服务器／存储／网络虚拟化、Hypervisor、容器为通行分类（教材未展开）\\[2pt]
\textbf{平台与安全：}AWS（2006 年推出，世界最大公有云平台服务提供商）· 阿里云（2009 年创立，中国最大）· 微软 Azure · Google 云 · 百度云 · 腾讯云；云安全与数据隐私风险
\end{knowmap}

\begin{keyterms}
\textbf{IaaS}：基础设施即服务，位于最底层，提供服务器、存储、网络等基本计算资源。\\
\textbf{PaaS}：平台即服务，位于第二层，提供分布式操作系统、分布式数据库等基础软件与开发环境。\\
\textbf{SaaS}：软件即服务，位于第三层，用户通过浏览器或终端直接使用远程服务器上的应用。\\
\textbf{虚拟化}：在软件中仿真计算机硬件、以虚拟资源为用户提供服务的计算形式。\\
\textbf{Hypervisor}：虚拟机监视器，运行在物理硬件与客户操作系统之间、负责创建调度和隔离虚拟机的软件层。\\
\textbf{虚拟机}：由虚拟化软件模拟出的、拥有独立操作系统和完整硬件视图的计算环境。\\
\textbf{容器}：共享宿主机操作系统内核、以进程级隔离运行应用的轻量级虚拟化形式。\\
\textbf{服务器虚拟化}：把一台物理服务器的资源抽象为多个相互隔离的虚拟机，提高资源利用率。\\
\textbf{存储虚拟化}：把多个异构存储设备抽象为统一的存储池，便于集中管理和在线扩容。\\
\textbf{网络虚拟化}：在物理网络之上构建相互隔离的虚拟网络，如虚拟私有云 VPC。\\
\textbf{AWS}：亚马逊网络服务，2006 年推出，教材称为世界上最大的公有云平台服务提供商。\\
\textbf{阿里云}：2009 年创立，中国最大的云计算平台，服务范围覆盖全球 200 多个国家和地区。\\
\textbf{云安全}：通过大量客户端监测软件行为异常，把木马病毒信息推送服务器端分析并分发解决方案。\\
\textbf{云存储}：通过集群应用、网格技术或分布式文件系统，把网络中大量存储设备集合起来协同提供存储与访问服务。
\end{keyterms}

\begin{contrast}{IaaS / PaaS / SaaS（每行依次为：所处层次·提供给用户的内容·典型例子）}
\citem{IaaS}{第 1 层（最底层）}{服务器、存储、网络等基本计算资源}{EC2 虚拟机实例、EBS 数据块存储}
\citem{PaaS}{第 2 层}{分布式操作系统、分布式数据库等基础软件与开发分发环境}{亚马逊平台服务（教材称为 PaaS 商业模式）}
\citem{SaaS}{第 3 层}{通过网页浏览器或终端使用的应用软件}{微信、微博、云存储、在线视频、Google Apps、Dropbox、iCloud}
\end{contrast}

\begin{contrast}{服务器虚拟化 / 存储虚拟化 / 网络虚拟化（每行依次为：虚拟对象·主要作用·典型形态；教材未展开，按通行分类整理）}
\citem{服务器虚拟化}{CPU、内存、磁盘、网络等物理服务器资源}{把一台服务器划分为多个虚拟机，提高利用率、隔离故障、快速部署}{VMware、KVM、Xen 等之上的多台虚拟机}
\citem{存储虚拟化}{多个异构存储设备}{抽象为统一存储池，集中管理、在线扩容、提高利用率}{分布式存储池、存储网关}
\citem{网络虚拟化}{物理网络链路与设备}{在同一物理网络上构建相互隔离的虚拟网络}{虚拟私有云 VPC、VLAN、软件定义网络}
\end{contrast}

\begin{contrast}{虚拟机 / 容器（每行依次为：隔离与资源单位·启动与开销·适用场景；教材未展开，按通行资料整理）}
\citem{虚拟机}{虚拟化硬件，每个虚拟机含完整客户操作系统}{启动较慢、镜像大、资源开销高、隔离性强}{运行不同操作系统的应用、强隔离多租户场景}
\citem{容器}{共享宿主机操作系统内核，以进程为隔离单位}{秒级启动、镜像小、开销低、隔离性相对弱}{微服务、持续交付、弹性伸缩的应用部署}
\end{contrast}

\qtypename{一、单项选择题}{10 题，每题 3 分，共 30 分}
\q 在 NIST 划分的云计算服务模式中，位于第 2 层、向用户提供分布式操作系统和分布式数据库等基础软件的是（\quad）
\choicesline{A. IaaS}{B. PaaS}{C. SaaS}{D. DaaS}

\q 用户购买服务器基础服务后即可自由安装和定制操作系统，不必理会硬件购买、托管、监控与维护，这种模式是（\quad）
\choicesline{A. IaaS}{B. PaaS}{C. SaaS}{D. 私有云}

\q 用户消费的服务完全从桌面系统或智能手机获取，如微信、微博、云存储、在线视频，这类服务属于（\quad）
\choicesline{A. IaaS}{B. PaaS}{C. SaaS}{D. 混合云}

\q 教材 4.1.6 中，亚马逊提供的 EC2 虚拟机实例属于（\quad）
\choicesline{A. IaaS}{B. PaaS}{C. SaaS}{D. 绿色节能技术}

\q 教材指出，因创立 PaaS 商业模式而被《商业周刊》评为 2008 年创新企业的是（\quad）
\choicesline{A. 亚马逊}{B. 微软}{C. Google}{D. 阿里云}

\q 在虚拟化技术中，Hypervisor（虚拟机监视器）位于（\quad）
\choices{A. 应用程序与数据库之间}{B. 物理硬件与客户操作系统之间}
{C. 浏览器与 Web 服务器之间}{D. 数据库与存储设备之间}

\q 容器与虚拟机的主要区别是（\quad）
\choices{A. 容器必须安装完整的客户操作系统}{B. 容器共享宿主操作系统内核，无需安装完整客户操作系统}
{C. 容器不能运行在 Linux 上}{D. 容器不属于虚拟化技术}

\q 教材指出，中国最大的云计算平台、创立于 2009 年的是（\quad）
\choicesline{A. 百度云}{B. 腾讯云}{C. 阿里云}{D. 华为云}

\q 教材称为世界上最大的公有云平台服务提供商的是（\quad）
\choicesline{A. AWS}{B. 微软 Azure}{C. Google 云}{D. 阿里云}

\q 教材 4.1.7 所述“云安全”的基本思路是（\quad）
\choices{A. 禁止用户把任何数据放到云端}{B. 通过大量客户端监测软件行为异常，服务器端分析后分发解决方案}
{C. 只依靠防火墙隔离内外网}{D. 把数据全部加密后离线保存}

\qtypename{二、多项选择题}{3 题，每题 4 分，共 12 分；每题有两个或两个以上正确项}
\q 关于 SaaS，下列说法正确的有（\quad）
\choices{A. 大多通过网页浏览器实现}{B. 用户消费的服务从桌面系统或智能手机中获取}
{C. 微信、微博、云存储属于 SaaS}{D. 用户必须自行维护服务器操作系统和数据库}

\q 关于虚拟化技术，下列说法正确的有（\quad）
\choices{A. 是在软件中仿真计算机硬件、以虚拟资源提供服务}{B. 可以把一台性能强大的服务器虚拟成多个独立小服务器}
{C. 可以把多台服务器虚拟成一台强大的服务器}{D. 虚拟化就是云计算}

\q 教材 4.1.6 中，属于 AWS 提供的服务有（\quad）
\choicesline{A. EC2 虚拟机实例}{B. EBS 数据块存储}{C. Elastic MapReduce 分析}{D. 微信}

\qtypename{三、判断题}{6 题，每题 2 分，共 12 分；正确的画“√”，错误的画“×”并写出改正后的表述}
\q PaaS 位于 IaaS 之上、SaaS 之下，向用户提供分布式操作系统、分布式数据库等基础软件。\quad（\quad）\quad 改正：\fixblank

\q 使用 IaaS 的用户必须自己购买、托管和维护物理服务器。\quad（\quad）\quad 改正：\fixblank

\q 容器与虚拟机一样，每个容器都需安装完整客户操作系统。\quad（\quad）\quad 改正：\fixblank

\q AWS 于 2006 年推出，教材称它是世界上最大的公有云平台服务提供商。\quad（\quad）\quad 改正：\fixblank

\q 阿里云创立于 2009 年，是中国最大的云计算平台，服务范围覆盖全球 200 多个国家和地区。\quad（\quad）\quad 改正：\fixblank

\q 云安全就是把数据放到云端，因此云上数据绝对安全。\quad（\quad）\quad 改正：\fixblank

\qtypename{四、填空题}{6 空，每空 1 分，共 6 分}
\q 按 NIST 划分，云计算的服务模式分为 IaaS、PaaS 和\fillblanks[2.0cm]三大类。

\q IaaS 的中文名称是\fillblanks[3.2cm]。

\q PaaS 的英文全称是\fillblanks[5.2cm]。

\q 在软件中仿真计算机硬件、以虚拟资源为用户提供服务的计算形式称为\fillblanks[2.0cm]；把一台物理服务器划分为多个虚拟服务器的做法属于\fillblanks[2.2cm]虚拟化。

\q 通过集群应用、网格技术或分布式文件系统，把网络中大量存储设备集合起来协同提供存储与访问服务的系统称为\fillblanks[2.2cm]。

\newpage
\qtypename{五、简答题}{4 题，每题 6 分，共 24 分}
\q 简述 IaaS、PaaS、SaaS 三层分别向用户提供什么、用户承担哪些职责，并各举一个教材中的例子。\anslines{2.4cm}

\q 什么是虚拟化？简述其两种应用模式与最大好处。\anslines{2.4cm}

\q 简述虚拟机与容器的主要区别（教材未展开，可按通行资料作答）。\anslines{2.4cm}

\q 简述教材 4.1.6 中 AWS 与阿里云的地位和主要服务，并列举教材提到的其他主流云平台。\anslines{2.4cm}

\newpage
\qtypename{六、综合题}{1 题，共 16 分}
\q 某农业企业拟把三套业务上云：（1）面向公众的农产品电商网站；（2）面向内部员工的 ERP 与办公系统；（3）需要弹性算力的育种数据分析平台。企业还计划把部分公开的农业技术资料做成在线课程。请回答：
（1）分别判断三套业务及在线课程最适合的服务模式（IaaS／PaaS／SaaS），并说明理由。（6 分）
（2）说明虚拟化在其中的作用；若企业自建私有云，服务器虚拟化能带来哪些好处？（4 分）
（3）从部署模式看这些业务应如何组合？并分析云安全与数据隐私风险及应对措施。（6 分）\anslines{5cm}

\end{document}
